\chapter{Techniques}

\section{Searching}
	\kactlimport{BinarySearchAnswer.cpp}
	\kactlimport{TernarySearch.cpp}

\section{Arrays}
	\kactlimport{PrefixSums.cpp}
	\kactlimport{TwoPointers.cpp}
	\textbf{Coordinate compression.} If values are up to $10^9$ but only their order matters, replace each by its
	index among the sorted distinct values: \texttt{vals = a; sort, erase(unique); x = lower\_bound(all(vals), x) - vals.begin()}.

	\textbf{Monotonic stack.} For the nearest bigger element to the left of each a[i], keep a stack of indices:
	pop while \texttt{a[top] <= a[i]}, then the top (if any) is the answer, then push i. $O(n)$ in total.

\section{Dynamic programming}
	\textbf{Designing a DP.}
	(1) State: the few facts about a partial solution that decide what can still happen,
	e.g.\ \texttt{dp[i][c]} = the best using the first $i$ items with capacity $c$.
	(2) Transition: get a state from smaller ones, usually by trying each choice for the next step.
	(3) Base cases. (4) An order in which every state comes after the states it uses, or recursion with
	memoisation, which finds the order by itself. (5) Where the answer is, e.g.\ \texttt{dp[n][W]}.
	The time is the number of states times the work per state: check it against the constraints.
	Counting problems add over the choices (mod $p$), optimisation problems take the min or max.

	Common states: a prefix, \texttt{dp[i]} (e.g.\ grid paths: \texttt{dp[i][j] = dp[i-1][j] + dp[i][j-1]};
	longest increasing subsequence: \texttt{LIS.h}); a node of a DAG, in topological order;
	a subset of up to 20 items, \texttt{dp[S]}; a segment, \texttt{dp[l][r]}. Games where both players play
	optimally: a position is winning iff some move leads to a losing position (usually, having no move loses).
	\kactlimport{Knapsack.cpp}
	\kactlimport{LCS.cpp}
	\kactlimport{LIS.h}
	\kactlimport{LDS.cpp}
	\kactlimport{BitmaskDP.cpp}
	\kactlimport{IntervalDP.cpp}
